\ficha{Gontijo (2014), Os mecanismos de funcionamento do padrão-ouro}{gontijo2014}

\begin{identificacao}
GONTIJO, Cláudio. Os mecanismos de funcionamento do ``padrão-ouro'': uma visão
crítica. \emph{Economia e Sociedade}, Campinas, v. 23, n. 1 (50), p. 243-280,
abr. 2014.

Artigo teórico e de história do pensamento econômico, 38 páginas, em português.
Lido integralmente na versão publicada.
\end{identificacao}

\bloco{Problema e tese}
O artigo pergunta como o padrão-ouro de fato ajustava o balanço de pagamentos, e
responde que não era pelo mecanismo de fluxo de metais e ajuste de preços
atribuído a Hume. Esse mecanismo supõe implicitamente moeda fiduciária e ignora a
lei do refluxo de Adam Smith, que governa qualquer regime de dinheiro-mercadoria
com moeda de crédito conversível. Onde o refluxo opera, a oferta de moeda é
endógena, o nível de preços fica ancorado no custo de produção do ouro, e o
ajuste externo se dá por fluxo de metal, manipulação da taxa de juros e dos
prazos do crédito, e entesouramento do público. A variação de preços existe, mas
é instrumental e se reverte ao final do processo.

\bloco{Argumento}
\begin{enumerate}[leftmargin=*,nosep]
\item Gontijo recusa a definição corrente do regime, que o faz nascer por volta
de 1880 com a adesão de vários países a três regras, conversibilidade a preço
fixo, liberdade de importar e exportar ouro, e vínculo entre circulação e estoque
metálico. Essa caracterização, diz ele, ``confunde dinheiro com crédito'' e
obscurece a lógica das crises periódicas do período (p. 246). Em lugar dela
define o padrão-ouro por três traços institucionais: dinheiro-mercadoria de livre
cunhagem, plena conversibilidade da moeda de crédito bancária, e liberdade de
exportar e importar metal (p. 246). O recorte do artigo começa em meados da
década de 1860, quando o cabo telegráfico e a integração do mercado de capitais
tornam eficiente a arbitragem de ouro e de juros.

\item Reconstitui o modelo tradicional já com a taxa de juros incorporada, adição
tornada necessária pela mobilidade do capital no fim do século XIX. Superávit
comercial traz ouro, os bancos ampliam prazos e cortam a taxa de desconto, sobem
demanda agregada, renda e preços, crescem as importações e caem as exportações
até o equilíbrio; no déficit ocorre o inverso, com corte de crédito e alta de
juros (p. 258).

\item Confronta esse modelo com a evidência que a própria literatura ortodoxa
reuniu. Triffin observa paralelismo, e não divergência, entre exportações e
importações dos vários países e entre suas taxas de desconto; os ajustes
descendentes de salários raramente foram amplos; as políticas de neutralização
antecedem a Primeira Guerra (p. 259). Whale e Bloomfield acrescentam que os
bancos centrais líderes não seguiam as chamadas regras do jogo, que o nível
nacional de preços era determinado pelo sistema mundial de preços, e que não há
evidência de efeito de longo prazo da disciplina monetária sobre o balanço de
pagamentos (p. 260).

\item Localiza a raiz do desacordo na teoria quantitativa embutida no mecanismo,
que supõe não haver distinção qualitativa entre moeda de ouro e moeda de crédito,
supõe a oferta de moeda independente da demanda, e supõe exógenas as velocidades
de circulação (p. 261). As três hipóteses caem. A moeda de crédito substitui o
dinheiro apenas na medida em que os pagamentos se compensam. A velocidade varia
com a \emph{echéance}, o prazo dos empréstimos, que os bancos alongam quando há
excesso de reservas e encurtam quando há aperto. E o crédito comercial, base do
crédito bancário via desconto de letras, é determinado pelo ritmo da produção,
isto é, pela demanda, o que explica cofres cheios e circulação reduzida nas
recessões do século XIX (p. 262).

\item O erro maior é ignorar a lei do refluxo, formulada por Smith e retomada por
Ricardo: os bancos ``nunca poderiam emitir mais notas do que o valor da moeda
metálica que teria circulado se não houvesse qualquer banco'' (p. 263), pois o
excesso volta ao emissor para conversão. Daí a conclusão que interessa aqui.

\chave{o ajustamento do balanço de pagamentos não se opera através da flutuação
do nível de preços, mas, através do fluxo de metal do país deficitário para o
superavitário, o que preserva o nível de preços}{p. 263}

\item Na prática, os bancos de emissão ajustam juros, prazos, limites
quantitativos e elegibilidade de ativos para preservar a conversibilidade, que
era o objetivo dominante. O próprio aumento da oferta de moeda desencadeia as
forças que a fazem refluir, de modo que a ``oferta de moeda'' era inteiramente
endógena (p. 266). O ajuste é real, mas não automático nem suave: se faz por
crises periódicas, conjugado ao ciclo de negócios. Legislação de lastro rígido é
dispensável e contraproducente, e o Bank Act de 1844, inspirado na currency
school, teve de ser suspenso em 1847, 1857 e 1866 (p. 266).
\end{enumerate}

\chave{não havia mudança permanente do nível de preços, que se mantinha estável e
vinculado ao custo relativo de produção do ouro ao nível mundial}{p. 273}

\bloco{Conceitos e definições operacionais}
Padrão-ouro: dinheiro-mercadoria de livre cunhagem sem senhoriagem, plena
conversibilidade das notas e dos créditos em conta corrente, e liberdade de
movimento do metal (p. 246). Moeda de crédito: promessas de pagamento, letras de
câmbio, notas conversíveis e depósitos, que substituem o dinheiro apenas na
medida em que os pagamentos se compensam. Lei do refluxo: emissão excedente
retorna ao banco para conversão e sai em remessa ao exterior, o que restabelece a
paridade sem alterar o nível de preços (p. 263). \emph{Echéance}: prazo do
desconto, instrumento de controle do crédito tão relevante quanto a taxa.
Endogeneidade: ``a oferta se ajustava à demanda'' (p. 264). Automatismo: o autor
recusa tanto a leitura inteiramente automática quanto a inteiramente gerenciada
(nota 22, p. 254).

\bloco{Evidência e método}
Não há dados novos, série reconstruída nem teste estatístico. O método combina
crítica textual dos clássicos, Hume, Smith, Ricardo, Stuart Mill e Marx, com uma
narrativa da história monetária britânica de 1821 a 1907 apoiada em fontes
secundárias clássicas, Feavearyear, King e Gregory, Scammell e Hawtrey. A
evidência empírica contra o modelo tradicional é toda de segunda mão, tomada de
Triffin, Bloomfield e Whale. Não há contrafactual.

\bloco{Agentes e vertentes}
De um lado, a tradição quantitativista e o mecanismo de Hume: Viner, o Cunliffe
Committee de 1918, Bordo, Eichengreen, e a currency school do Peel Act. De outro,
a linhagem que Gontijo assume, Smith, Ricardo e Marx, com Glasner e Laidler na
leitura contemporânea da endogeneidade. Bloomfield, Triffin e Whale entram como
ortodoxos cuja evidência desmente o modelo que eles próprios descrevem. Em nota,
opõe Mundell e Laffer, partidários do automatismo, a Hawtrey e Bloomfield,
partidários do sistema gerenciado (p. 254). Lindert é citado pela assimetria
entre Londres, que financiava o comércio exterior em libras e gerenciava o
sistema, e os demais centros, com poder bem menor de atrair capital por elevação
de juros (p. 254).

\bloco{Críticas e limites}
O artigo é sobre o centro. Toda a evidência histórica vem do Reino Unido, e não há
uma linha sobre a periferia, sobre a América Latina ou sobre o Brasil. Não trata
de credibilidade, reputação, risco soberano ou prêmio de juros, eixos de Bordo e
Rockoff e de Obstfeld e Taylor, nem da assimetria centro-periferia de Franco
(1988): a assimetria que reconhece é a de Lindert, entre praças financeiras. A lei
do refluxo, como formulada, pressupõe conversibilidade plena nos dois sentidos e
liberdade de exportação de metal. Um arranjo de conversibilidade limitada, com
teto de emissão e taxa fixada por lei abaixo do par legal, não é o caso do
modelo, e transferir a conclusão para ele exige argumento próprio. Por fim, as
afirmações sobre velocidade de circulação e \emph{echéance} apoiam-se em relato
narrativo, não em medida.

\begin{dialogo}
\item Elasticidade da circulação. Procurar nos jornais o argumento de que a
emissão da Caixa contra ouro não pode ser inflacionária porque o papel excedente
volta ao emissor, ou porque a circulação se ajusta às necessidades do comércio.
Vocabulário de época: ``necessidades do comércio'', ``meio circulante'', ``papel
lastrado''. Testa se a defesa da Caixa mobiliza raciocínio de tipo refluxo, com
oferta endógena, ou se aceita a premissa quantitativista dos adversários.

\item Leitura da entrada e da saída de ouro. Nos boletins e comentários sobre o
movimento diário da Caixa, verificar se a saída de ouro é lida como sinal
automático de excesso de emissão, na chave de Hume, ou como efeito do ciclo, da
safra e da conta de capital. Testa a hipótese de que o mecanismo de Hume era o
senso comum dos publicistas de 1906.

\item Juros e crédito bancário. Procurar menção a desconto, prazos, redesconto,
aperto do mercado monetário e ao papel do Banco do Brasil ao lado do argumento
cambial. Se o canal do crédito aparece, o debate reconhecia o mecanismo que
Gontijo identifica como o efetivo; se não aparece, o horizonte da discussão era
estritamente cambial.

\item Lastro integral e limite de emissão. Procurar o argumento de que só o
lastro metálico integral impede o abuso da emissão. Gontijo trata essa posição
como currency school e a considera contraproducente, lembrando as três suspensões
do Bank Act de 1844. Testa se a defesa brasileira do lastro reproduz a doutrina
inglesa e com que fundamento.

\item Crise e ciclo internacional. Verificar se a crise de 1907 e a suspensão de
1914 são atribuídas pelos jornais a causas domésticas ou a um ciclo mundial.
Testa, no material brasileiro, a tese do paralelismo entre países que Gontijo
toma de Triffin.
\end{dialogo}

\usonocapitulo{Parte I. O texto não trata do Brasil e não sustenta nenhuma
afirmação sobre o caso brasileiro. Serve para duas coisas. Primeira, mostrar que
o mecanismo de Hume, pressuposto implícito de boa parte da literatura da Parte I
e do vocabulário de época, é contestado dentro da própria teoria monetária, o que
autoriza tratar os argumentos de 1906 sobre elasticidade da circulação como
posições teóricas disputáveis e não como erro de leigo. Segunda, fornecer o
vocabulário analítico, refluxo, endogeneidade da oferta de moeda e prazo do
crédito, com que ler as afirmações de época sobre emissão contra ouro, entrada e
saída de metal e crédito bancário. Funciona como contrapeso interno à Parte I
diante de Bordo e Kydland e de Bordo e Rockoff, que assumem o regime como regra
crível sem examinar o mecanismo de ajuste.}
